\chapter{VAR e Causalidade de Granger}
\label{cap:var}
% ── índice ──────────────────────────────────────────────────────
\index{VAR|textbf}
\index{vetor autoregressivo|see{VAR}}
\index{var@\texttt{var()}|textbf}
\index{causalidade de Granger|textbf}
\index{granger@\texttt{granger()}|textbf}
\index{função impulso-resposta|textbf}
\index{irf@\texttt{irf()}|textbf}
\index{decomposição da variância}

% ─────────────────────────────────────────────────────────────────────────────
\section{O modelo VAR}

O modelo Autorregressivo Vetorial (\textbf{VAR($p$)}) generaliza o AR para
sistemas multivariados. Com $K$ variáveis e $p$ defasagens:

\[
  \mathbf{y}_t = \mathbf{c} + \mathbf{A}_1 \mathbf{y}_{t-1}
                 + \cdots + \mathbf{A}_p \mathbf{y}_{t-p}
                 + \boldsymbol{\varepsilon}_t
\]

onde $\mathbf{y}_t = (y_{1t}, \ldots, y_{Kt})'$ é o vetor de variáveis,
$\mathbf{A}_j$ são matrizes $K \times K$ de coeficientes e
$\boldsymbol{\varepsilon}_t \sim (0, \Sigma_u)$ é o vetor de choques não
correlacionados no tempo mas potencialmente correlacionados entre si.

O VAR é estimado equação por equação via OLS — cada equação é um AR($p$)
nos $K$ defasados. A condição de estacionariedade exige que todas as raízes
características de $\mathbf{A}_1 + \cdots + \mathbf{A}_p$ estejam dentro do
círculo unitário.

% ─────────────────────────────────────────────────────────────────────────────
\section{Causalidade de Granger}

$x$ \textbf{causa Granger} $y$ se as defasagens de $x$ têm poder preditivo
para $y$ além do que as defasagens de $y$ sozinhas fornecem. Formalmente,
testa-se a hipótese:
\[
  H_0: A_{12}(1) = A_{12}(2) = \cdots = A_{12}(p) = 0
\]
na equação de $y_1$, onde $A_{12}(j)$ é o coeficiente de $y_{2,t-j}$.

O teste é um teste F padrão na equação auxiliar:
\[
  F = \frac{(RSS_r - RSS_u)/p}{RSS_u / (T - 2p - 1)} \sim F(p,\; T-2p-1)
\]

\begin{nota}
  Causalidade de Granger é \emph{previsibilidade}, não causalidade
  estrutural. $x$ que antecede $y$ e correlaciona com ela passa no teste,
  mesmo sem relação causal direta.
\end{nota}

% ─────────────────────────────────────────────────────────────────────────────
\section{Função de Resposta ao Impulso}

A \textbf{IRF} (\emph{Impulse Response Function}) rastreia o efeito de um
choque unitário em $y_k$ sobre todas as variáveis ao longo do tempo.

A representação MA($\infty$) do VAR é $\mathbf{y}_t = \sum_{h=0}^\infty
\mathbf{\Phi}_h \boldsymbol{\varepsilon}_{t-h}$, com $\mathbf{\Phi}_0 = I$
e $\mathbf{\Phi}_h = \mathbf{A}_1 \mathbf{\Phi}_{h-1} + \cdots +
\mathbf{A}_p \mathbf{\Phi}_{h-p}$. A IRF ortogonalizada usa a decomposição
de Cholesky de $\Sigma_u$ para isolar choques não correlacionados — a
\textbf{ordenação de Cholesky} (ordem das variáveis) importa.

% ─────────────────────────────────────────────────────────────────────────────
\section{Verificação por DGP}

O DGP usa um VAR(1) bivariado com estrutura de causalidade assimétrica:
$y_2$ causa $y_1$ (coeficiente $0{,}2$), mas $y_1$ causa $y_2$
fracamente (coeficiente $0{,}1$):

\[
  \begin{pmatrix} y_{1t} \\ y_{2t} \end{pmatrix}
  = \begin{pmatrix} 0{,}5 & 0{,}2 \\ 0{,}1 & 0{,}6 \end{pmatrix}
    \begin{pmatrix} y_{1,t-1} \\ y_{2,t-1} \end{pmatrix}
  + \begin{pmatrix} \varepsilon_{1t} \\ \varepsilon_{2t} \end{pmatrix}
\]

\begin{lstlisting}[caption={DGP VAR(1) — Granger e IRF}]
seed(42)
input df
id
1
end
for i in 1..=9 { df = append(df, df) }
df |> filter(_n <= 300)
generate df e1 = rnormal()
generate df e2 = rnormal()
generate df y1 = 0.0
generate df y2 = 0.0
generate df t  = _n
let n = nrow(df)
let i = 2
while i <= n {
    let l1 = mean(df, y1, if = t == i - 1)
    let l2 = mean(df, y2, if = t == i - 1)
    let a  = mean(df, e1, if = t == i)
    let b  = mean(df, e2, if = t == i)
    replace df y1 = 0.5*l1 + 0.2*l2 + a if t == i
    replace df y2 = 0.1*l1 + 0.6*l2 + b if t == i
    i = i + 1
}
let m = var(df, y1, y2, lags=1)
display m
granger(df, y1, y2, lags=1)
granger(df, y2, y1, lags=1)
irf(m, steps=6)
\end{lstlisting}

\subsection*{Ajuste do modelo}

\begin{lstlisting}[language={}, basicstyle=\ttfamily\small,
  frame=none, numbers=none, backgroundcolor=\color{codbg},
  xleftmargin=0pt, xrightmargin=0pt, breaklines=false]
VAR(1)  N=299  AIC=-5.0269  BIC=-4.9526

Sigma_u:
  [ 0.0819  0.0046 ]
  [ 0.0046  0.0772 ]
\end{lstlisting}

A matriz de covariância residual $\hat\Sigma_u$ é quase diagonal — os
choques $\varepsilon_1$ e $\varepsilon_2$ são praticamente ortogonais,
como especificado no DGP.

\subsection*{Causalidade de Granger}

\begin{lstlisting}[language={}, basicstyle=\ttfamily\small,
  frame=none, numbers=none, backgroundcolor=\color{codbg},
  xleftmargin=0pt, xrightmargin=0pt, breaklines=false]
================== Granger Causality Test ==================
  H₀: y2 não causa Granger y1   (lags=1)
  F(1, 296) = 9.4437   p = 0.0023
  Conclusion: REJEITA H₀ — y2 causa Granger y1

================== Granger Causality Test ==================
  H₀: y1 não causa Granger y2   (lags=1)
  F(1, 296) = 4.5768   p = 0.0332
  Conclusion: REJEITA H₀ — y1 causa Granger y2
\end{lstlisting}

Ambas as direções rejeitam $H_0$: $y_2 \to y_1$ com $p = 0{,}002$ e
$y_1 \to y_2$ com $p = 0{,}033$. O DGP impõe $\phi_{12} = 0{,}2$ e
$\phi_{21} = 0{,}1$ — a causalidade bidirecional é detectada, embora
$y_2 \to y_1$ seja mais forte.

\subsection*{Resposta ao impulso}

\begin{lstlisting}[language={}, basicstyle=\ttfamily\small,
  frame=none, numbers=none, backgroundcolor=\color{codbg},
  xleftmargin=0pt, xrightmargin=0pt, breaklines=false]
IRF — VAR(1) — 6 passos

  Impulso: y1
  ──────────────────────────────────────
       h          y1          y2
  ──────────────────────────────────────
       1      0.2862      0.0161
       2      0.1430      0.0382
       3      0.0758      0.0364
       4      0.0426      0.0286
       5      0.0251      0.0207
       6      0.0154      0.0145
  ──────────────────────────────────────

  Impulso: y2
  ──────────────────────────────────────
       h          y1          y2
  ──────────────────────────────────────
       1      0.0000      0.2774
       2      0.0405      0.1595
       3      0.0432      0.0958
       4      0.0352      0.0595
       5      0.0260      0.0378
       6      0.0183      0.0243
  ──────────────────────────────────────
\end{lstlisting}

Um choque em $y_1$ tem efeito rápido sobre $y_1$ (decai de $0{,}29$ para
$0{,}015$ em 6 períodos) e efeito pequeno sobre $y_2$. Um choque em $y_2$
transmite-se progressivamente a $y_1$ — atinge $0{,}043$ no período 3,
depois decai — capturando o canal $\phi_{12} = 0{,}2$ do DGP.

% ─────────────────────────────────────────────────────────────────────────────
\section{Seleção de defasagens}

\begin{lstlisting}[caption={Comparação por número de defasagens}]
let m1 = var(df, y1, y2, lags=1)
let m2 = var(df, y1, y2, lags=2)
let m3 = var(df, y1, y2, lags=3)
esttab(m1, m2, m3)
\end{lstlisting}

Seleciona-se a defasagem com menor BIC (parcimônia) ou AIC (ajuste).
Com o DGP VAR(1), espera-se que $p=1$ minimize ambos.

% ─────────────────────────────────────────────────────────────────────────────
\section{Sintaxe}

\begin{lstlisting}
var(df, y1, y2, lags=1)
var(df, y1, y2, y3, lags=2)
granger(df, y, x, lags=1)
irf(m, steps=10)
\end{lstlisting}

\begin{lstlisting}[caption={Fluxo típico — PIB e inflação}]
load "macro.csv" as df

adf(df, pib)
adf(df, ipca)

let m = var(df, pib, ipca, lags=2)
granger(df, pib, ipca, lags=2)
granger(df, ipca, pib, lags=2)
irf(m, steps=12)
\end{lstlisting}

\begin{nota}
  O VAR exige séries \emph{estacionárias}. Se as variáveis forem I(1) e
  cointegradas, o modelo adequado é o VECM (cap.~\ref{cap:coint}). Se
  forem I(1) mas não cointegradas, diferencie antes de estimar.
\end{nota}

\section{Validação empírica}

O DGP VAR(1) deste capítulo corresponde ao caso de validação
\texttt{var\_book} em \texttt{validation/cases/}. O caso compara a estimação
OLS equação por equação do \hayashi{} com referências em R e Python. Execute
\lstinline{hay validate} para ver o estado atual.
